\documentclass[11pt]{article}
\usepackage[margin=1in]{geometry}
\usepackage{amsmath,amssymb,amsthm}
\usepackage{booktabs}
\usepackage{graphicx}
\usepackage{microtype}
\usepackage[hidelinks]{hyperref}
\usepackage{titlesec}
\usepackage{caption}

\captionsetup{font=small,labelfont=bf}
\titlespacing*{\section}{0pt}{1.1em}{0.5em}
\titlespacing*{\subsection}{0pt}{0.8em}{0.35em}
\setlength{\abovedisplayskip}{5pt}
\setlength{\belowdisplayskip}{5pt}
\setlength{\parskip}{0.25em}

\newtheorem{proposition}{Proposition}

\title{\vspace{-2.5em}\bfseries A Telegraph-Equation Price Model with Stochastic\\
Variance and Double-Exponential Jumps:\\
Specification, Validation, and a Structural Null Result}

\author{David Belyn\thanks{Correspondence: \texttt{dbelyn@hotmail.com}. Working paper, draft.}}

\date{September 2026}

\begin{document}
\maketitle
\vspace{-1.8em}

\begin{abstract}
\noindent
We specify and validate the Dirac--Heston--Jump (DHJ) model, a hybrid
transition-density model for foreign exchange that replaces the
Black--Scholes diffusion with a Goldstein--Kac telegraph process on a
log-price lattice, couples it to a Heston stochastic-variance path, and
overlays Kou double-exponential jumps. The telegraph representation gives
finite propagation speed and non-Gaussian tails at finite relaxation rate
$\kappa$, and recovers Black--Scholes exactly as $\kappa \to \infty$. We give
the full continuum and discrete specification, a Monte-Carlo-outer /
PDE-inner numerical scheme, and a seven-part validation suite: the model
reproduces Garman--Kohlhagen call prices to $0.15\%$ relative error at
$N=1600$ lattice sites, satisfies the martingale condition
$\mathbb{E}[S_T]=S_0 e^{(r_d-r_f)T}$ to within $0.003\%$, preserves that
condition under compensated jumps to $0.07\%$, and exhibits zero negative
densities and zero measurable boundary mass loss.

The model was then deployed as a directional forecaster and produced
$49$--$50\%$ accuracy over more than $2{,}500$ logged one-day forecasts. We
show this was not a calibration failure but structurally predetermined, and
prove it two ways. First, we derive in closed form that the model's
``chiral charge'' signal satisfies
$\bar{Q}_5(T) = \delta_{cp}\,e^{-(2\kappa+r_d)T}$ exactly: it is the initial
spinor asymmetry times a deterministic scalar, carrying no information from
the variance process, the jumps, or the PDE itself. Since $\delta_{cp}$ was
seeded from RSI, the apparatus computed a damped momentum oscillator; we
confirm the identity numerically to six significant figures. Second, a
correctly specified risk-neutral model is a martingale by construction, so
its directional content is exactly the carry; at a one-day horizon this caps
achievable accuracy at $50.28\%$ against a break-even threshold of $51.08\%$
after spread, an expectancy of $-0.89$ pips per trade. The validation suite
that certified the model as correct is the same evidence proving it cannot
be directionally profitable at that horizon. We argue this is a general trap
in applying derivative-pricing machinery to directional forecasting.
\end{abstract}

\vspace{-0.8em}
\section{Introduction}

Option-pricing models are calibrated to reproduce the risk-neutral
transition density $p^{\mathbb{Q}}(S_T \mid S_0)$. Directional trading
requires the real-world density $p^{\mathbb{P}}$. The two coincide only up
to a market price of risk that pricing models, by design, do not estimate.
This paper documents what happens when a carefully validated $\mathbb{Q}$-measure
model is deployed as a $\mathbb{P}$-measure forecaster, and extracts the
general lesson.

The model studied here, Dirac--Heston--Jump (DHJ), was built to address
three documented deficiencies of Black--Scholes in FX: Gaussian tails,
constant volatility, and continuous paths. The first is addressed by
replacing the diffusion with a telegraph (Goldstein--Kac) process, which
propagates at finite speed $c$ and relaxes to diffusion at rate $2\kappa$,
producing excess kurtosis and skew of order $1/\kappa$; the second by a
Heston CIR variance path; the third by a Kou jump overlay. Each ingredient
is standard. The contribution is the combination, the numerical treatment
required to make it stable, and the negative result with its structural
explanation.

\section{Model Specification}
\label{sec:spec}

\subsection{State variables}

Let $x = \log(S/S_0)$ be the log-price relative to spot and $\tau$ calendar
time in years. The state is a two-component real field
$\psi = (\psi_+, \psi_-)$ on $x$, where $\psi_+$ is right-moving probability
flux and $\psi_-$ is left-moving flux. Both are real and non-negative: this
is a Euclidean (imaginary-time) formulation, so the components are genuine
densities, not quantum amplitudes. We define
\begin{equation}
P(x,\tau) = \psi_+ + \psi_-, \qquad
Q_5(x,\tau) = \psi_+ - \psi_-,
\end{equation}
the total density and the \emph{chiral density}. The aggregate chiral charge
is $\bar{Q}_5(\tau) = \int Q_5(x,\tau)\,dx$.

\subsection{Telegraph dynamics}

The field evolves under the coupled system
\begin{align}
\partial_\tau \psi_+ &= D\,\partial_x^2 \psi_+ - c\,\partial_x \psi_+
  - \nu\,\partial_x \psi_+ - \kappa \psi_+ + \gamma \psi_- - r_d \psi_+,
  \label{eq:psip}\\
\partial_\tau \psi_- &= D\,\partial_x^2 \psi_- + c\,\partial_x \psi_-
  - \nu\,\partial_x \psi_- - \kappa \psi_- + \gamma \psi_+ - r_d \psi_-,
  \label{eq:psim}
\end{align}
with $D$ the diffusion coefficient, $c$ the propagation speed, $\kappa$ the
flip (relaxation) rate, $\gamma$ the chiral mixing rate, $\nu$ the log-price
drift, and $r_d$ the domestic rate entering as discounting. Setting
$\gamma = \kappa$ conserves total probability up to discounting. In
$(P, Q_5)$ variables, \eqref{eq:psip}--\eqref{eq:psim} become
\begin{align}
\partial_\tau P &= D\,\partial_x^2 P - c\,\partial_x Q_5 - \nu\,\partial_x P - r_d P,
  \label{eq:Peq}\\
\partial_\tau Q_5 &= D\,\partial_x^2 Q_5 - c\,\partial_x P - \nu\,\partial_x Q_5
  - 2\kappa Q_5 - r_d Q_5.
  \label{eq:Qeq}
\end{align}
For $D = \nu = r_d = 0$, eliminating $Q_5$ gives the classical telegraph
equation
\begin{equation}
\partial_\tau^2 P + 2\kappa\,\partial_\tau P = c^2\,\partial_x^2 P,
\end{equation}
which interpolates between ballistic propagation at speed $c$ (small
$\kappa$) and diffusion (large $\kappa$).

\subsection{The Black--Scholes limit and the finite-$\kappa$ correction}

As $\kappa \to \infty$ the chiral mode relaxes instantaneously. Setting
$\partial_\tau Q_5 \approx 0$ and dropping terms small against $2\kappa Q_5$
in \eqref{eq:Qeq} gives $Q_5 \approx -(c/2\kappa)\,\partial_x P$.
Substituting into \eqref{eq:Peq},
\begin{equation}
\partial_\tau P = \Big(\underbrace{D + \tfrac{c^2}{2\kappa}}_{D_{\text{eff}}}\Big)\partial_x^2 P
  - \nu\,\partial_x P - r_d P.
\label{eq:bslimit}
\end{equation}
With $D = \sigma^2/2$ and $c = \sigma$ this is the exact Black--Scholes
Fokker--Planck equation plus an excess-diffusion term. Since $c^2 = 2D$,
\begin{equation}
D_{\text{eff}} = D\big(1 + \kappa^{-1}\big),
\qquad
\sigma_{\text{eff}}^2 = \sigma^2\big(1 + \kappa^{-1}\big),
\end{equation}
so the finite-mass model carries a relative variance excess of exactly
$1/\kappa$, vanishing in the diffusive limit. Skewness enters at order
$1/\kappa$ through the asymmetric momentum sector. These are the intended
non-Gaussian corrections.

\subsection{Heston stochastic variance}

The instantaneous variance $v(\tau)$ follows a CIR process,
\begin{equation}
dv = \kappa_H(\theta - v)\,d\tau + \xi \sqrt{v}\,dW_v,
\end{equation}
with mean-reversion speed $\kappa_H$, long-run variance $\theta$,
vol-of-vol $\xi$, and $dW_v$ correlated with the price shock $dW_S$ at
$\rho = \mathrm{corr}(dW_S, dW_v)$. It is integrated by a Milstein scheme
with full truncation,
\begin{equation}
v_{n+1} = \max\!\Big(v_n + \kappa_H(\theta - v_n)\,\Delta\tau
  + \xi\sqrt{v_n^+}\,\Delta W
  + \tfrac{1}{4}\xi^2\big[(\Delta W)^2 - \Delta\tau\big],\ 0\Big),
\end{equation}
where $v_n^+ = \max(v_n, 0)$ and $\Delta W = z_n\sqrt{\Delta\tau}$ with
$z_n \sim \mathcal{N}(0,1)$. Milstein gives strong order $1.0$ against
Euler's $0.5$.

\subsection{Variance decomposition: avoiding a double count}
\label{sec:decomp}

The Heston price shock decomposes into correlated and independent
components,
\begin{equation}
dW_S = \rho\,dW_v + \sqrt{1-\rho^2}\,dW_\perp .
\end{equation}
The correlated part is applied separately as a path-conditional
translation (Section~\ref{sec:lev}), so the PDE must evolve only the
independent part. Using the full $v$ would double-count the variance
$\rho^2 v$. The PDE coefficients are therefore
\begin{equation}
D(\tau) = \frac{(1-\rho^2)\,v(\tau)}{2},
\qquad
c(\tau) = \sqrt{(1-\rho^2)\,v(\tau)} .
\label{eq:Dc}
\end{equation}

\subsection{Regime-dependent mass}

The relaxation rate is made variance-dependent so that calm markets are
more ballistic (trending, low $\kappa$) and stressed markets more diffusive,
\begin{equation}
\kappa(v) = \min\!\Big(\kappa_0 + \frac{\kappa_1}{v},\ \frac{1}{2\,\Delta\tau}\Big),
\qquad \gamma = \kappa .
\label{eq:kappa}
\end{equation}
The $\kappa_1/v$ term is singular as $v \to 0$; the cap is required for
stability (Section~\ref{sec:stab}).

\subsection{Kou jumps and the compensator}

Jumps arrive as a Poisson process of intensity $\lambda$ per year. Each jump
$Y$ in log-price has the double-exponential density
\begin{equation}
f_Y(y) = p\,\eta_+ e^{-\eta_+ y}\,\mathbf{1}_{\{y>0\}}
       + (1-p)\,\eta_- e^{\eta_- y}\,\mathbf{1}_{\{y<0\}},
\qquad \eta_+ > 1,\ \eta_- > 0 .
\end{equation}
Mean jump sizes are $1/\eta_+$ up and $1/\eta_-$ down; asymmetry
$\eta_- < \eta_+$ encodes crash skew. The martingale compensator is
\begin{equation}
\zeta \;=\; \mathbb{E}\big[e^{Y} - 1\big]
\;=\; \frac{p\,\eta_+}{\eta_+ - 1} + \frac{(1-p)\,\eta_-}{\eta_- + 1} - 1,
\label{eq:comp}
\end{equation}
and $\lambda\zeta$ is subtracted from the drift. The condition $\eta_+ > 1$
is necessary for $\mathbb{E}[e^Y]$ to exist and is enforced at construction.

\subsection{Drift and the Garman--Kohlhagen boundary}

Under the FX risk-neutral measure with domestic rate $r_d$ and foreign rate
$r_f$, the carry is $r_d - r_f$ and the log-price drift is
\begin{equation}
\nu(\tau) = (r_d - r_f) - \tfrac{1}{2}v(\tau) - \lambda\zeta ,
\label{eq:nu}
\end{equation}
which enforces $\mathbb{E}[S_T] = S_0\,e^{(r_d - r_f)T}$. Setting $r_f = 0$
reduces the model to the equity convention. The reference density and call
price against which the model is validated are Garman--Kohlhagen,
\begin{equation}
p_{\mathrm{GK}}(S) = \frac{1}{S\sigma\sqrt{2\pi T}}
  \exp\!\left[-\frac{\big(\log(S/S_0) - (r_d - r_f - \tfrac{1}{2}\sigma^2)T\big)^2}
  {2\sigma^2 T}\right],
\end{equation}
\begin{equation}
C_{\mathrm{GK}} = S_0 e^{-r_f T}\Phi(d_1) - K e^{-r_d T}\Phi(d_2),
\qquad
d_{1,2} = \frac{\log(S_0/K) + \big(r_d - r_f \pm \tfrac{1}{2}\sigma^2\big)T}
  {\sigma\sqrt{T}} .
\end{equation}

\subsection{Nested limits}

The specification nests three known models exactly:
$\kappa_0 \to \infty,\ \kappa_1 = \xi = \lambda = \rho = 0$ gives
Black--Scholes; $\kappa_0 \to \infty,\ \kappa_1 = \lambda = 0$ gives Heston;
$\xi = \lambda = \rho = 0$ gives the deterministic telegraph model. These
limits are the basis of the validation suite.

\section{Numerical Scheme}
\label{sec:num}

The solver is Monte Carlo on the outer loop (variance paths) and PDE on the
inner loop (the spinor field conditional on a variance path).

\subsection{Lattice and initial condition}

The domain is $x \in [-x_h, x_h]$ with
$x_h = w\sigma_0\sqrt{T} + \tfrac{1}{2}$, $w$ half-widths in standard
deviations (default $6$), discretised at $N$ sites with spacing
$\Delta x = 2x_h/(N-1)$ and Dirichlet boundaries. The initial condition is a
Gaussian of width $\Delta x$ centred at $x = 0$, split asymmetrically,
\begin{equation}
\psi_\pm(x, 0) \;\propto\; \theta_{\pm}\,
  \exp\!\left(-\frac{x^2}{2\Delta x^2}\right),
\qquad
\theta_+ = \theta_{ic},\quad \theta_- = 1 - \theta_{ic},
\end{equation}
\begin{equation}
\theta_{ic} = \operatorname{clip}\!\big(\tfrac{1}{2} + \tfrac{1}{2}\delta_{cp},\ 0.01,\ 0.99\big),
\label{eq:thetaic}
\end{equation}
normalised so $\int P\,dx = 1$. A one-cell width is the narrowest smearing
that yields a smooth density while keeping the artificial initial variance
$O(\Delta x^2)$ far below $\sigma^2 T$; a width of two or more cells injects
$3$--$7\%$ excess variance at short horizons.

\subsection{Spatial discretisation and integrator}

With $f_j = f(x_j)$, the operators are
\begin{equation}
\Delta^2 f_j = \frac{f_{j+1} - 2f_j + f_{j-1}}{\Delta x^2},
\qquad
\nabla^0 f_j = \frac{f_{j+1} - f_{j-1}}{2\Delta x} .
\end{equation}
Both the wave term and the drift term use the second-order central
difference $\nabla^0$. First-order upwinding would be stable but introduces
numerical diffusion $D_{\text{num}} = c\,\Delta x/2$, which at realistic
$\Delta x$ is comparable to the physical $D$ and would corrupt the very
non-Gaussian corrections the model exists to produce. The Wilson diffusion
term already damps the high-wavenumber modes that make central differencing
unstable in pure advection. Time integration is classical fourth-order
Runge--Kutta on the semi-discrete system.

\subsection{Stability}
\label{sec:stab}

Three constraints are enforced. The CFL condition for the explicit scheme
is
\begin{equation}
\Delta\tau \;\le\; 0.45\,\frac{\Delta x^2}{2D + c\,\Delta x} .
\label{eq:cfl}
\end{equation}
The global step is set from a stationary worst-case variance. The CIR
stationary distribution is Gamma with standard deviation
$\sigma_v = \xi\sqrt{\theta/(2\kappa_H)}$, so we take
$v_{\text{ref}} = \max(\theta + 2\sigma_v,\ v_0,\ \theta)$ and size
$\Delta\tau$ from \eqref{eq:cfl} at $v_{\text{ref}}$. When a realised path
excursion exceeds $v_{\text{ref}}$, the step is subdivided adaptively into
$n_{\text{sub}} = \lceil \Delta\tau / \Delta\tau_{\text{CFL}}(v_i)\rceil$
substeps, capped at $32$.

Second, the reaction term $-2\kappa$ is stiff, with RK4 absolute-stability
limit $|2\kappa h| < 2.785$. The cap in \eqref{eq:kappa} at
$\kappa \le 1/(2\Delta\tau)$ gives $|2\kappa h| \le 1$, a $2.8\times$
margin, while still permitting the large $\kappa_0$ needed for
Black--Scholes recovery tests.

Third, the drift is clamped at $|\nu| \le \Delta x/(2\Delta\tau)$ as a
sanity bound; with \eqref{eq:nu} it is never binding in practice.

\subsection{Leverage as a grid translation}
\label{sec:lev}

The correlated shock contributes a path-conditional drift
$\rho\sqrt{v}\,dW_v$. Discretising it as a PDE advection term would impose a
CFL constraint scaling with $|\rho|\sqrt{v}$, which at realistic FX
parameters dominates \eqref{eq:cfl} and forces a much smaller step. Instead
it is applied exactly, as a translation of the field after each step:
\begin{equation}
\Delta x_{\text{lev}} = \rho\,\sqrt{v_i}\;z_i\,\sqrt{\Delta\tau},
\end{equation}
implemented by sub-cell linear interpolation. Writing
$\Delta x_{\text{lev}}/\Delta x = n_s + \phi$ with $n_s = \lfloor\cdot\rfloor$
and $\phi \in [0,1)$,
\begin{equation}
\tilde{\psi}_j = (1-\phi)\,\psi_{j - n_s} + \phi\,\psi_{j - n_s - 1} .
\label{eq:shift}
\end{equation}
Translation is unconditionally stable and mathematically equivalent to the
advection term, removing the constraint entirely. Kou jumps are applied by
the same operator, sampling $n \sim \text{Poisson}(\lambda\Delta\tau)$ jumps
per step and translating by each $Y$.

\subsection{Monte Carlo accumulation}

For each of $N_p$ variance paths the field is evolved to $T$ and accumulated
with weight equal to its retained norm $\|\psi^{(k)}\|_1$, rather than
normalising each path first:
\begin{equation}
\bar{\psi}_j = \frac{\sum_{k} \psi^{(k)}_j}{\sum_{k}\|\psi^{(k)}\|_1},
\qquad \|\psi^{(k)}\|_1 = \textstyle\sum_j \big(\psi^{(k)}_{+,j} + \psi^{(k)}_{-,j}\big)\Delta x .
\end{equation}
Per-path normalisation divides by a near-zero norm on paths that have leaked
off the domain, amplifying numerical noise catastrophically;
weight-based accumulation lets such paths contribute in proportion to the
mass they retain. Since total mass decays as $e^{-r_d T}$ under discounting,
boundary leakage is diagnosed as
\begin{equation}
\text{mass loss} = \max\!\Big(0,\ 1 - \frac{\int \bar{P}\,dx}{e^{-r_d T}}\Big).
\end{equation}

\section{Validation}
\label{sec:val}

All results below were regenerated from the current build. Test~1 fixes
$\xi = \lambda = \rho = 0$ and $\kappa_0 = 500$, where the model must
reproduce Garman--Kohlhagen exactly, and refines the lattice
(Table~\ref{tab:conv}). The relative call-price error falls by roughly a
factor of three per doubling of $N$ over the first two refinements and
reaches $0.15\%$ at $N = 1600$.

\begin{table}[ht]
\centering\small
\begin{tabular}{rccccc}
\toprule
$N$ & $\mathbb{E}[S_T]$ & GK forward & $C_{\text{DHJ}}$ & $C_{\text{GK}}$ & rel.\ err.\\
\midrule
200  & 1.001063 & 1.001042 & 0.010286 & 0.009935 & 3.54\,\% \\
400  & 1.001048 & 1.001042 & 0.010027 & 0.009935 & 0.93\,\% \\
800  & 1.001044 & 1.001042 & 0.009965 & 0.009935 & 0.30\,\% \\
1600 & 1.001043 & 1.001042 & 0.009950 & 0.009935 & 0.15\,\% \\
\bottomrule
\end{tabular}
\caption{Garman--Kohlhagen recovery against lattice resolution.
$S_0=1$, $\sigma=0.082$, $r_d=0.0525$, $r_f=0.0400$, $T=1/12$.}
\label{tab:conv}
\end{table}

Test~2 sweeps $\kappa_0 \in \{10, \dots, 1000\}$ at $N=400$ and confirms the
limit \eqref{eq:bslimit}: relative error falls monotonically
$3.29\% \to 1.68\% \to 1.29\% \to 1.07\% \to 0.93\% \to 0.89\%$, with zero
negative density cells throughout. The residual at large $\kappa_0$ is the
$N=400$ discretisation floor from Table~\ref{tab:conv}, not a limit failure.

Test~3 checks the martingale condition across horizons and vol-of-vol.
Relative deviations of $\mathbb{E}[S_T]$ from $S_0 e^{(r_d-r_f)T}$ are
$0.003\%$ ($T=30$d, $\xi=0$), $0.002\%$ ($30$d, $\xi=0.3$), $0.016\%$
($91$d, $\xi=0.3$) and $0.001\%$ ($30$d, $\kappa_0=50$). Test~4 repeats this
with jumps active at $\lambda \in \{1,3,5,10\}$, verifying the compensator
\eqref{eq:comp}: errors are $0.015\%$, $0.071\%$, $0.014\%$, $0.019\%$,
all within Monte Carlo error at $N_p = 300$.

Test~5 finds zero negative density cells across eight configurations
spanning $N \in [200, 800]$, $T \in \{1, 3\}$ months, high $\xi$ and
$\lambda$, and $\rho = -0.6$. Test~6 measures boundary mass loss at
$0.0000\%$ for all four standard configurations. The model is, as a
numerical implementation of its own specification, correct.

\section{Signal Extraction, and Where It Fails}
\label{sec:sig}

\subsection{The extraction layer}

Three quantities were read from the solved density. The probability of
finishing above spot,
\begin{equation}
\Pi_{\uparrow} \;=\; \int_{S_0}^{\infty} p_{\mathrm{DHJ}}(S)\,dS
\;\approx\; \sum_{j\,:\,S_j \ge S_0} p_j\,(S_{j+1} - S_j),
\end{equation}
the expected move $\mathbb{E}[S_T] - S_0$, and the aggregate chiral charge
$\bar{Q}_5(T)$, interpreted as a momentum or sentiment proxy. The initial
asymmetry was seeded from the $14$-period RSI with a MACD agreement filter,
\begin{equation}
\delta_{cp} = \operatorname{clip}\!\left(\frac{\mathrm{RSI} - 50}{100}
  \cdot \big(1 - \tfrac{1}{2}\mathbf{1}_{\{\mathrm{sgn}(\mathrm{RSI}-50)\,\neq\,\mathrm{sgn}(\mathrm{MACD}_h)\}}\big),
  \ -\tfrac{1}{2},\ \tfrac{1}{2}\right),
\label{eq:deltacp}
\end{equation}
and signals were assigned by thresholds on $\bar{Q}_5$ with confirmation
from $\Pi_\uparrow$: strong signals at $|\bar{Q}_5| > 0.15$ with
$\Pi_\uparrow$ confirming beyond $0.52$ or $0.48$, mild signals at
$|\bar{Q}_5| > 0.05$, otherwise neutral.

\subsection{The chiral charge is not a model output}

\begin{proposition}
\label{prop:q5}
Under Dirichlet boundaries and $\gamma = \kappa$, the aggregate chiral
charge obeys, exactly,
\begin{equation}
\bar{Q}_5(T) \;=\; \delta_{cp}\;e^{-(2\kappa + r_d)T}.
\label{eq:q5closed}
\end{equation}
\end{proposition}

\begin{proof}
Integrate \eqref{eq:Qeq} over $x$. The terms
$\int \partial_x^2 Q_5\,dx$, $\int \partial_x P\,dx$ and
$\int \partial_x Q_5\,dx$ are boundary terms and vanish under Dirichlet
conditions, leaving
$\tfrac{d}{d\tau}\bar{Q}_5 = -(2\kappa + r_d)\,\bar{Q}_5$.
The leverage and jump operators are translations \eqref{eq:shift}, which
preserve the integral, so they do not enter. At $\tau = 0$, by
\eqref{eq:thetaic}, $\bar{Q}_5(0) = \theta_+ - \theta_- = 2\theta_{ic} - 1
= \delta_{cp}$. Integrating gives \eqref{eq:q5closed}.
\end{proof}

We verified \eqref{eq:q5closed} against the compiled engine at
$\delta_{cp} \in \{-0.5, \dots, 0.5\}$, $\sigma = 0.082$, $T = 1/365$,
$\kappa_0 = 1.0$, $\kappa_1 = 0.002$, giving
$\kappa = 1.2974$ and a predicted factor $e^{-(2\kappa+r_d)T} = 0.992773$.
The ratio of measured to predicted charge was $1.00000$ at every point, to
all six digits reported.

The consequence is severe. The chiral charge, which determined the
direction of every signal the model emitted, is the RSI-derived seed
multiplied by a deterministic scalar within $0.8\%$ of unity at a one-day
horizon. It contains no information from the Heston path, the jump process,
the leverage coupling, or the PDE evolution. Inverting the thresholds
through \eqref{eq:deltacp}, the signal rule was exactly
\begin{equation}
\text{STRONG} \iff |\mathrm{RSI} - 50| > 15.11,
\qquad
\text{MILD} \iff |\mathrm{RSI} - 50| > 5.04,
\end{equation}
doubling to $30.2$ and $10.1$ when MACD disagreed. A lattice PDE with
stochastic variance and compensated jumps was computing
$\mathrm{RSI} > 65$.

This is a specification error of a particular and instructive kind. Nothing
in Section~\ref{sec:val} could have caught it: the model was correct, the
solver was correct, and the quantity being read out was simply not a
function of either. The conserved-quantity analysis that proves
\eqref{eq:q5closed} takes one page and was not performed until after
deployment.

\section{Empirical Results}
\label{sec:emp}

The model ran in production against a live FX feed, logging a forecast per
instrument per hourly cycle with a one-day evaluation horizon, scored on
realised direction. Accuracy was assessed with the Wilson score lower bound
at $z = 1.64$,
\begin{equation}
p_{\text{LB}} = \frac{\hat{p} + \frac{z^2}{2n}
  - z\sqrt{\frac{\hat{p}(1-\hat{p})}{n} + \frac{z^2}{4n^2}}}
  {1 + \frac{z^2}{n}},
\end{equation}
which is preferred to the normal interval because it does not degenerate at
small $n$ or extreme $\hat p$: a cell at $4/5$ carries $p_{\text{LB}} \approx 0.38$
and is correctly not mistaken for an edge.

Over more than $2{,}500$ evaluated one-day forecasts the model achieved
$49$--$50\%$ directional accuracy, indistinguishable from a coin flip. A
conditional analysis partitioned forecasts into cells by signal strength,
volatility regime, session and instrument. Several cells showed in-sample
accuracy of $66$--$77\%$; on out-of-sample data the same cells realised
$15$--$31\%$. The sign reversal, rather than mere decay toward $50\%$, is
the signature of fitting noise across a large cell count, and is what
motivated adopting pre-registered interpretation rules for all subsequent
tests. An independent five-voter ensemble built to replace the model scored
$49.8\%$ on $n = 325$ out-of-sample forecasts, $p_{\text{LB}} = 0.453$.

As a positive control, the same evaluation harness applied to index
buy-and-hold recovered a Sharpe ratio of $0.55$, confirming that the
detector was capable of finding a real effect where one exists.

\section{Why the Null Result Was Structural}
\label{sec:why}

Two arguments, one specific and one general, show the outcome was not a
calibration failure.

The specific argument is Proposition~\ref{prop:q5}: the directional signal
was an RSI threshold rule, and momentum rules of that form in liquid FX have
been repeatedly documented to carry no exploitable edge net of costs.

The general argument is that the model was, by construction, a martingale.
Test~3 confirms $\mathbb{E}[S_T] = S_0 e^{(r_d - r_f)T}$ to $0.003\%$. Its
entire directional content is therefore the carry, and nothing else. Under
the model's own log-normal $\mathbb{Q}$-density, the implied probability of
an up move is
\begin{equation}
\Pi_\uparrow^{\mathbb{Q}} = \Phi\!\left(\frac{m}{s}\right),
\qquad
m = \big(r_d - r_f - \tfrac{1}{2}\sigma^2\big)T,
\qquad
s = \sigma\sqrt{T}.
\end{equation}
Against a symmetric entry with a one-standard-deviation target and stop
$d = S_0\sigma\sqrt{T}$ and round-trip cost $C$, expectancy is
$d(2\Pi_\uparrow - 1) - C$, so the break-even accuracy is
\begin{equation}
\Pi^{*} = \tfrac{1}{2} + \frac{C}{2d}.
\end{equation}
Table~\ref{tab:horizon} evaluates both for EUR/USD at $\sigma = 0.082$,
$r_d - r_f = 125$\,bp, $S_0 = 1.08$, $C = 1.2$ pips.

\begin{table}[ht]
\centering\small
\begin{tabular}{lccrr}
\toprule
Horizon & $\Pi_\uparrow^{\mathbb{Q}}$ & $\Pi^{*}$ (break-even)
  & edge (pp) & $\mathbb{E}$[pips] \\
\midrule
1 day    & 0.5028 & 0.5108 & $-0.80$ & $-0.89$ \\
1 week   & 0.5063 & 0.5048 & $+0.15$ & $+0.36$ \\
1 month  & 0.5128 & 0.5023 & $+1.05$ & $+5.36$ \\
3 months & 0.5222 & 0.5014 & $+2.09$ & $+18.48$ \\
1 year   & 0.5444 & 0.5007 & $+4.37$ & $+77.38$ \\
\bottomrule
\end{tabular}
\caption{Achievable directional accuracy under the model's own risk-neutral
density versus the accuracy required to break even after spread.}
\label{tab:horizon}
\end{table}

At the one-day horizon at which the model was deployed, its own maximum
attainable accuracy of $50.28\%$ falls $0.80$ percentage points short of the
$51.08\%$ required to cover the spread, an expectancy of $-0.89$ pips per
trade. The observed $49$--$50\%$ is fully consistent with this bound. No
improvement in calibration, lattice resolution or jump parameters could have
changed it, because the deficit is a property of the risk-neutral
specification rather than of its numerical solution. The result was
derivable in closed form before a single forecast was logged, from
quantities the validation suite was already reporting.

Table~\ref{tab:horizon} also shows the diagnosis. The gross edge turns
positive beyond roughly one week and grows with horizon, because the carry
signal accumulates linearly in $T$ while noise grows as $\sqrt{T}$ and the
spread is paid once. The model was not wrong; it was deployed at the one
horizon where its signal is structurally smaller than its transaction cost.
That positive gross expectancy should not be read as free money: at longer
horizons this is simply the carry trade, whose returns are widely understood
as compensation for crash risk. Our own backtests over 2014--2026 found
time-series carry at Sharpe $0.00$ and cross-sectional carry at
$0.14 \pm 0.26$, consistent with a risk premium that is real but not large
relative to its variance.

\section{Discussion}

Three lessons generalise beyond this model.

\emph{Validation proves consistency, not usefulness.} Every test in
Section~\ref{sec:val} asks whether the implementation matches the
specification. None asks whether the specification answers the question
being posed. A model can pass a complete numerical validation suite and
still be structurally incapable of the task it was built for, and ours did.

\emph{Risk-neutral machinery does not transfer to directional forecasting.}
The martingale property is the defining virtue of a pricing model and a
disqualifying defect in a forecaster. Practitioners moving pricing code into
a signal pipeline should state the measure change explicitly and estimate
it, rather than reading $\mathbb{P}$-measure conclusions off a
$\mathbb{Q}$-measure density.

\emph{Conserved quantities should be checked before deployment, not after.}
Proposition~\ref{prop:q5} is elementary. Had the integral of \eqref{eq:Qeq}
been taken at specification time, the entire signal layer would have been
recognised as a reparameterised RSI before any code was written.

The model retains genuine value in the role it was designed for. It produces
a validated, arbitrage-consistent transition density with tunable tails and
skew, it prices European options against a Garman--Kohlhagen benchmark to
sub-percent accuracy, and the parameters $\kappa$, $\xi$, $\rho$ and
$\lambda$ give independent control over excess kurtosis and skew for
volatility-surface work. The natural next application is variance
forecasting rather than direction: returns in liquid FX appear
unforecastable at short horizons, but realised variance is strongly
persistent, and the model's variance layer is the component the validation
suite most directly certified.

\begin{thebibliography}{99}
\small
\setlength{\itemsep}{0pt}\setlength{\parsep}{0pt}\setlength{\parskip}{0pt}
\bibitem{bailey2014} Bailey, D.~H. and M.~L\'opez de Prado (2014).
  The deflated Sharpe ratio. \emph{Journal of Portfolio Management} 40(5), 94--107.
\bibitem{bns2004} Barndorff-Nielsen, O.~E. and N.~Shephard (2004).
  Power and bipower variation with stochastic volatility and jumps.
  \emph{Journal of Financial Econometrics} 2(1), 1--37.
\bibitem{bs1973} Black, F. and M.~Scholes (1973). The pricing of options and
  corporate liabilities. \emph{Journal of Political Economy} 81(3), 637--654.
\bibitem{corsi2009} Corsi, F. (2009). A simple approximate long-memory model
  of realized volatility. \emph{Journal of Financial Econometrics} 7(2), 174--196.
\bibitem{fama1970} Fama, E.~F. (1970). Efficient capital markets.
  \emph{Journal of Finance} 25(2), 383--417.
\bibitem{gk1983} Garman, M.~B. and S.~W.~Kohlhagen (1983). Foreign currency
  option values. \emph{Journal of International Money and Finance} 2(3), 231--237.
\bibitem{goldstein1951} Goldstein, S. (1951). On diffusion by discontinuous
  movements and on the telegraph equation. \emph{Quarterly Journal of
  Mechanics and Applied Mathematics} 4(2), 129--156.
\bibitem{harvey2016} Harvey, C.~R., Y.~Liu and H.~Zhu (2016).
  $\ldots$ and the cross-section of expected returns.
  \emph{Review of Financial Studies} 29(1), 5--68.
\bibitem{heston1993} Heston, S.~L. (1993). A closed-form solution for options
  with stochastic volatility. \emph{Review of Financial Studies} 6(2), 327--343.
\bibitem{kac1974} Kac, M. (1974). A stochastic model related to the telegrapher's
  equation. \emph{Rocky Mountain Journal of Mathematics} 4(3), 497--509.
\bibitem{kou2002} Kou, S.~G. (2002). A jump-diffusion model for option pricing.
  \emph{Management Science} 48(8), 1086--1101.
\bibitem{lord2010} Lord, R., R.~Koekkoek and D.~van Dijk (2010). A comparison of
  biased simulation schemes for stochastic volatility models.
  \emph{Quantitative Finance} 10(2), 177--194.
\bibitem{meese1983} Meese, R.~A. and K.~Rogoff (1983). Empirical exchange rate
  models of the seventies. \emph{Journal of International Economics} 14(1--2), 3--24.
\bibitem{milstein1975} Milstein, G.~N. (1975). Approximate integration of
  stochastic differential equations. \emph{Theory of Probability and its
  Applications} 19(3), 557--562.
\bibitem{stw1999} Sullivan, R., A.~Timmermann and H.~White (1999). Data-snooping,
  technical trading rule performance, and the bootstrap.
  \emph{Journal of Finance} 54(5), 1647--1691.
\bibitem{white2000} White, H. (2000). A reality check for data snooping.
  \emph{Econometrica} 68(5), 1097--1126.
\bibitem{wilson1927} Wilson, E.~B. (1927). Probable inference, the law of
  succession, and statistical inference. \emph{Journal of the American
  Statistical Association} 22(158), 209--212.
\bibitem{wilson1974} Wilson, K.~G. (1974). Confinement of quarks.
  \emph{Physical Review D} 10(8), 2445--2459.
\end{thebibliography}

\end{document}
